\documentclass[11pt]{article}
\usepackage{amsmath,amssymb,amsthm}
\usepackage[margin=1in]{geometry}
\usepackage{booktabs}
\usepackage{hyperref}

\newtheorem{theorem}{Theorem}
\newtheorem{finding}{Finding}
\newtheorem{nonclaim}{Nonclaim}

\title{Step 245: Independent Codification Pattern}
\author{anti\_loc framework track}
\date{2026-05-16}

\begin{document}
\maketitle

\section{Candidate Pattern}

\begin{finding}[Independent Codification Pattern, candidate]
Let \(N \ge 2\) typed cascades on distinct mathematical tracks independently codify analogs of the same structural typed condition in their own internal records before cross-track validation. Then the condition has convergent evidence as a framework-internal discipline rather than as a track-specific artifact.
\end{finding}

The finding is not a formal theorem of typed cascades. It is a corpus-pending recognition rule: independent local codification plus later cross-track validation is stronger evidence than cross-track validation alone.

\section{Primary Evidence: Bridge Impossibility}

\begin{center}
\begin{tabular}{p{0.14\linewidth}p{0.30\linewidth}p{0.46\linewidth}}
\toprule
Track & Local record & Independent codification \\
\midrule
RH & Step 189 / \texttt{findings\_rh.md} & A non-CRE carrier plus a typed bridge to Riemann \(\Xi_{BC}=0\) makes the composite CRE; the bridge is RH-strength. \\
Hodge & H-18R/H-3R/H-8R source-status rules & \texttt{known\_zero} is licensed only by a source-stated cycle-span theorem \(A_{X,p}=V_{X,p}\); nearby Hodge-theoretic facts do not promote without the target-strength bridge. \\
P-vs-NP & NP47 universal P-machine atlas & The universal P-machine atlas is target-equivalent: if SAT is in \(P\), one quotient is the SAT boundary; proving all such quotients miss the boundary is P-vs-NP-strength. \\
\bottomrule
\end{tabular}
\end{center}

These three codifications predate the cross-track Bridge Impossibility audits. They express the same typed condition in track-native vocabulary:
\[
  \text{non-target-equivalent closure} + \text{bridge to target}
  \quad\Rightarrow\quad
  \text{bridge/composite carries target strength}.
\]

\section{Other Typed Conditions}

CTMT recursion and CRCFT modes are verified cross-track, but the available records do not show comparable independent pre-codification in all tracks. Selberg-Class Dichotomy Generalization is an L-function-family pattern, not a cross-track independent codification. Cascade Reduction remains primarily RH-originated in the available records.

\section{Status}

\[
\boxed{\text{candidate (3 independent Bridge Impossibility codifications) corpus-pending}}
\]

\begin{nonclaim}
This step does not solve any target conjecture and does not assert universality of the pattern beyond the stated evidence. The strongest evidence is for Bridge Impossibility.
\end{nonclaim}

\section{Verdict}

\[
\texttt{V\_independent\_codification\_typed\_3\_instance\_BI}.
\]

\end{document}
